\documentclass[a4paper,11pt]{article}
\usepackage[utf8]{inputenc}
\usepackage[T1]{fontenc}
\usepackage[french]{babel}
\usepackage[left=2cm,right=2cm,top=2cm,bottom=2cm]{geometry}
\usepackage{xcolor}
\usepackage{hyperref}
\usepackage{fontawesome5}
\usepackage{parskip}

% --- Couleurs Airbus ---
\definecolor{primary}{RGB}{0, 32, 91}

\begin{document}

% --- EN-TÊTE ---
\begin{center}
    {\Large \textbf{Loïc Aron MBASSI EWOLO}}\\[4pt]
    \small
    \faMapMarker\ Cergy, France (Mobile France entière) \quad
    \faPhone\ (+33) 7 59 52 33 98 \quad
    \faEnvelope\ \href{mailto:loicaron.mbassiewolo@ensea.fr}{loicaron.mbassiewolo@ensea.fr}
\end{center}
\vspace{0.5cm}

% --- DESTINATAIRE ---
\hfill
\begin{minipage}[t]{0.45\textwidth}
    \textbf{Airbus Protect}\\
    Département Managed Services\\
    Toulouse, France
\end{minipage}

\vspace{0.8cm}

\textbf{Objet :} Candidature Stage Recherche IA / NLP -- LLM \& RAG (6 mois)\\
\textit{Réf : Stage Recherche en IA / NLP -- Cybersécurité}

\vspace{0.5cm}

Madame, Monsieur,

Votre offre de stage sur les mécanismes de génération avancés pour LLM, notamment les architectures RAG et la réduction des hallucinations, correspond exactement à mes travaux récents. Actuellement en double diplôme ingénieur à l'\textbf{ENSEA} et à l'\textbf{École Polytechnique de Yaoundé}, je suis disponible dès \textbf{avril 2026} pour 6 mois.

J'ai eu la chance de travailler au \textbf{MILA} (Institut Québécois d'IA) en tant qu'assistant de recherche. J'y ai étudié le phénomène de Grokking (généralisation tardive des réseaux de neurones), reproduit des expériences SOTA sous \textbf{PyTorch} et mené une veille bibliographique intensive. Cette expérience m'a appris la rigueur de la recherche : formulation d'hypothèses, expérimentation systématique et analyse critique des résultats.

En parallèle, j'ai développé des compétences concrètes en \textbf{NLP et LLM}. Mon projet de système multi-agents agricole utilise \textbf{RAG} pour enrichir les réponses de 5 agents orchestrés via Google ADK et Gemini. J'ai dû gérer les problématiques de cohérence des réponses et de résolution de conflits entre agents -- des défis proches de la vérifiabilité des contenus générés que vous mentionnez. Lors de la compétition MLPC, j'ai fine-tuné des modèles \textbf{BERT} et \textbf{GPT-2} pour détecter des contenus sexistes, ce qui m'a sensibilisé aux biais et aux limites des LLM.

Le contexte \textbf{cybersécurité} de votre stage m'intéresse particulièrement. L'étude des attaques contre les systèmes RAG (injection de logs malicieux, empoisonnement de bases de connaissances) représente un sujet fascinant où mes compétences en IA et ma curiosité pour la sécurité informatique se rejoignent.

Je maîtrise les outils que vous utilisez : \textbf{PyTorch}, \textbf{Hugging Face}, Python (NumPy, scikit-learn), et je suis à l'aise avec la lecture et la synthèse de papers en anglais. Mon expérience en animation de workshops ML/LLM à Polytechnique Yaoundé m'a également appris à communiquer sur des sujets techniques complexes.

Je serais ravi de discuter de vos travaux sur la génération hybride et adaptative lors d'un entretien.

Je vous prie d'agréer, Madame, Monsieur, l'expression de mes salutations distinguées.

\vspace{0.8cm}

\textbf{Loïc Aron MBASSI EWOLO}\\
\textit{Élève-Ingénieur ENSEA / Polytechnique Yaoundé}

\end{document}
